%! Function Model Selection and Construction — Unit 1, Lesson 1.7 — Guided Notes (Answer Key)
\documentclass[10pt]{article}
\usepackage{precalculus-article}
\usepackage{precalculus-key}   % pulls in -boxes; do NOT also load -boxes
\newcommand{\TallMath}[1]{$\displaystyle #1\rule[-1.4em]{0pt}{3.2em}$}

\begin{document}

\pageheader{Unit 1, Lesson 1.7}{Guided Notes --- Answer Key}

\vspace{0.12in}

\begin{objectivebox}
By the end of this lesson, I will be able to\ldots
\begin{enumerate}[itemsep=2pt]
  \item Decide what kind of function fits a table by taking \ans{first} and
        \ans{second} differences of \ans{equally} spaced inputs.
  \item Choose a model type from the \ans{context} alone --- area, volume,
        inverse proportion, or a rule that changes at a threshold.
  \item Build a model by writing the \ans{constraint}, solving it for one variable,
        and \ans{substituting} into the quantity I want.
  \item Use a model: evaluate it, find an average rate of change with
        \ans{units}, and state the \ans{assumption} it rests on.
\end{enumerate}
\end{objectivebox}

\vspace{0.1in}

\boxguard
\begin{vocabbox}
Fill in each term as we define it together.
\par\vspace{2pt}
\noindent\textbf{\textcolor{plum}{Function model:}}\\[1pt]
\ansline{A function built to stand in for a real situation, so it can be predicted and measured.}\\[3pt]
\noindent\textbf{\textcolor{plum}{First differences:}}\\[1pt]
\ansline{Differences between consecutive outputs of equally spaced inputs. Constant $\Rightarrow$ linear.}\\[3pt]
\noindent\textbf{\textcolor{plum}{Second differences:}}\\[1pt]
\ansline{Differences of the first differences. Constant $\Rightarrow$ quadratic.}\\[3pt]
\noindent\textbf{\textcolor{plum}{Constraint:}}\\[1pt]
\ansline{The equation the context forces on the variables --- it collapses two unknowns into one.}\\[3pt]
\noindent\textbf{\textcolor{plum}{Domain restriction:}}\\[1pt]
\ansline{The inputs the context allows, usually narrower than what the algebra allows.}\\[3pt]
\noindent\textbf{\textcolor{plum}{Model assumption:}}\\[1pt]
\ansline{A simplifying claim about the world that must hold for the model to be trusted.}
\end{vocabbox}

\vspace{0.1in}

\boxguard
\begin{hookbox}
The garden club has exactly \textbf{80 feet of fencing} for a rectangular pen.
One whole side runs along the barn wall, so that side needs no fence at all.

\smallskip
Before any algebra --- should the pen be a square? Vote, then write down
\textit{why} you voted that way.

\ansline{Most classes vote ``square,'' because a square is the best answer when all four sides are fenced.}
\ansline{Here one side is free, so the pen should run long along the barn. Revisited in Part 4.}
\end{hookbox}

\vspace{0.1in}

\boxguard[26]
\begin{notesbox}{1. Reading the Model Off a Table}
\textbf{The requirement, first.} Differences only mean something when the inputs
are \ans{equally} spaced. Check the $x$ row before you subtract anything.

\medskip
\textbf{Table A.} Take the first differences.

\smallskip
\renewcommand{\arraystretch}{1.4}
\begin{center}
\begin{tabular}{|c|c|c|c|c|c|}
\hline
$x$ & 0 & 1 & 2 & 3 & 4 \\
\hline
$y$ & 5 & 8 & 11 & 14 & 17 \\
\hline
1st difference & --- & \ans{3} & \ans{3} & \ans{3} & \ans{3} \\
\hline
\end{tabular}
\end{center}

\smallskip
The first differences are constant, so Table A is \ans{linear}.

\medskip
\textbf{Table B --- the garden club's five trial layouts.} Here $w$ is the width
in feet and $A$ is the area in square feet.

\smallskip
\renewcommand{\arraystretch}{1.4}
\begin{center}
\begin{tabular}{|c|c|c|c|c|c|}
\hline
$w$ (ft) & 5 & 10 & 15 & 20 & 25 \\
\hline
$A$ (ft$^2$) & 350 & 600 & 750 & 800 & 750 \\
\hline
1st difference & --- & \ans{250} & \ans{150} & \ans{50} & \ans{$-50$} \\
\hline
2nd difference & --- & --- & \ans{$-100$} & \ans{$-100$} & \ans{$-100$} \\
\hline
\end{tabular}
\end{center}

\smallskip
The first differences are \textbf{not} constant --- so we take differences again.
The second differences are all \ans{$-100$}, so Table B is \ans{quadratic}.

\medskip
\textbf{The decision tree.}
\begin{itemize}[itemsep=4pt]
  \item 1st differences constant $\Rightarrow$ \ans{linear} \quad (degree \ans{1})
  \item 2nd differences constant $\Rightarrow$ \ans{quadratic} \quad (degree \ans{2})
  \item 3rd differences constant $\Rightarrow$ \ans{cubic} \quad (degree \ans{3})
  \item $n$th differences constant $\Rightarrow$ a degree-\ans{$n$} polynomial
\end{itemize}

\medskip
\textbf{What it means.} Constant \textit{first} differences say the
\ans{amount} of change repeats. Constant \textit{second} differences say the
\ans{rate} of change is itself changing at a steady rate.
\end{notesbox}

\vspace{0.1in}

\boxguard[26]
\begin{notesbox}{2. Reading the Model Off the Context}
No table? The description carries the clue.

\smallskip
\renewcommand{\arraystretch}{1.5}
\begin{center}
\begin{tabularx}{0.94\linewidth}{@{} >{\raggedright\arraybackslash}p{0.52\linewidth} X @{}}
\hline
\textbf{What the context says} & \textbf{Model type} \\
\hline
The same amount is added each step & \ans{Linear} \\
\hline
An \textbf{area}, or two lengths multiplied & \ans{Quadratic} \\
\hline
A \textbf{volume}, or three lengths multiplied & \ans{Cubic} \\
\hline
Several turns --- it rises, falls, rises again & \ans{Higher-degree polynomial} \\
\hline
One quantity is \textbf{inversely proportional} to another & \ans{Rational} \\
\hline
The rule \textbf{changes at a threshold} & \ans{Piecewise} \\
\hline
\end{tabularx}
\end{center}

\medskip
\textbf{Inverse proportion} means $y = \dfrac{k}{x^n}$: as $x$ grows, $y$
\ans{shrinks}. These are \textbf{rational} functions --- Unit 3 is about them.
Today you only have to recognize the words.

\medskip
\textbf{A counting fact.} A degree-$n$ polynomial can be fit through
\ans{$n+1$} points with distinct inputs. So four points can always be fit by a
cubic --- which is exactly why ``it fits the data'' is not by itself a reason to
believe a model.
\end{notesbox}

\vspace{0.1in}

\boxguard[30]
\begin{notesbox}{3. Building the Model from the Constraint}
Back to the 80 feet of fencing. Let $w$ be the width and $\ell$ the length.

\smallskip
\begin{center}
\begin{tikzpicture}[x=1cm, y=1cm]
  \fill[lilac] (0,0) rectangle (5,1.9);
  \draw[very thick, slate] (-0.5,1.9) -- (5.5,1.9);
  \node[slate, above, font=\small] at (2.5,1.95) {barn wall --- no fencing needed};
  \draw[very thick, plum] (0,1.9) -- (0,0) -- (5,0) -- (5,1.9);
  \node[plum, left,  font=\small] at (0,0.95) {$w$};
  \node[plum, right, font=\small] at (5,0.95) {$w$};
  \node[plum, below, font=\small] at (2.5,0) {$\ell$};
\end{tikzpicture}
\end{center}

\medskip
\textbf{Step 1 --- write the constraint.} Two widths and one length use all the
fencing:
\begin{center}
\ans{2} $w \; + \;$ \ans{$\ell$} $= 80$
\end{center}

\textbf{Step 2 --- solve the constraint for one variable.}
\begin{center}
$\ell = $ \ans{$80 - 2w$}
\end{center}

\textbf{Step 3 --- substitute into the quantity you actually want.} We want
\textbf{area}, and area is $A = w \cdot \ell$.

\begin{work}
  A(w) &= w \cdot \ell \\
       &= w(80 - 2w) \\
       &= 80w - 2w^2
\end{work}

\begin{center}
$A(w) = $ \ans{$80w - 2w^2$}
\end{center}

\textbf{Step 4 --- restrict the domain from the context.} The width must be
positive, and the length $80 - 2w$ must be positive too:

\begin{center}
\ans{0} $< w <$ \ans{40}
\end{center}

Why the upper end? At that width there is \ans{no fencing left for the length}.

\medskip
\textbf{Now check it against Table B.} From the table, $A(20) = 800$. From the
model:

\begin{work}
  A(20) &= 80(20) - 2(20)^2 \\
        &= 1600 - 800 \\
        &= 800
\end{work}

\smallskip
The table and the constraint produced the \ans{same} function --- and Part 2
predicted its type, because area is \ans{two} lengths multiplied.
\end{notesbox}

\vspace{0.1in}

\boxguard[30]
\begin{notesbox}{4. Using the Model}
A model that is never used is just algebra. Ours is $A(w) = 80w - 2w^2$ on
$0 < w < 40$.

\medskip
\textbf{(a) Evaluate.}

\begin{work}
  A(10) &= 80(10) - 2(10)^2 \\
        &= 800 - 200 \\
        &= 600 \\
  A(30) &= 80(30) - 2(30)^2 \\
        &= 2400 - 1800 \\
        &= 600
\end{work}

$A(10) = $ \ans{600} ft$^2$ \qquad $A(30) = $ \ans{600} ft$^2$

\medskip
\textbf{(b) Average rate of change --- with units.}

\begin{work}
  \frac{A(20) - A(10)}{20 - 10} &= \frac{800 - 600}{10} \\
                                &= 20 \\
  \frac{A(30) - A(20)}{30 - 20} &= \frac{600 - 800}{10} \\
                                &= -20
\end{work}

On $[10, 20]$: \ans{20} \qquad On $[20, 30]$: \ans{$-20$}

\smallskip
Units: \ans{square feet of area per foot of width}

\smallskip
In words: between $w = 10$ and $w = 20$, each extra foot of width adds about
\ans{20} square feet of area --- but past $w = 20$ each extra foot
\ans{takes away from} the area.

\medskip
\textbf{(c) The sign change locates the maximum.} The rate went from
\ans{$+$} to \ans{$-$}, so the largest area happens \textbf{between} $w = 10$
and $w = 30$. For $A(w) = aw^2 + bw$ the vertex sits at $w = -\dfrac{b}{2a}$:

\begin{work}
  w &= -\frac{b}{2a} \\
    &= -\frac{80}{2(-2)} \\
    &= 20
\end{work}

Best width $w = $ \ans{20} ft, \quad length $\ell = 80 - 2(20) = $
\ans{40} ft, \quad area $= $ \ans{800} ft$^2$.

\smallskip
\textbf{Back to the vote.} The best pen is \ans{20} ft by \ans{40} ft
--- it is \textbf{not} a square, because one side is free.
\end{notesbox}

\vspace{0.1in}

\boxguard[26]
\begin{notesbox}{5. The Fine Print: Restrictions and Assumptions}
Every model carries restrictions the algebra will not tell you about.

\medskip
\begin{itemize}[itemsep=5pt]
  \item \textbf{Domain restriction:} $w$ is a width in feet, so
        \ans{0} $< w <$ \ans{40}. Values outside this are not bad
        arithmetic --- they are outside the \ans{domain}.
  \item \textbf{Range restriction:} the area satisfies \ans{0} $< A \le$
        \ans{800} square feet.
\end{itemize}

\medskip
\textbf{The assumption sentence.} Fill in the frame:

\smallskip
\begin{center}
\fbox{\begin{minipage}{0.92\linewidth}
\medskip
``This model assumes that \ans{all 80 ft is used} , and it is only trustworthy for
inputs \ans{between 0 and 40 ft} .''
\medskip
\end{minipage}}
\end{center}

\medskip
Name two things the pen model assumes about the real world:

\ansline{Every foot of fence is used; the pen is exactly rectangular, with no gate gap.}

\medskip
\textbf{When one rule is not enough.} Suppose the barn wall is only 30 feet
long. Then a pen with $\ell > 30$ needs fencing on part of the fourth side too,
and the area rule \textbf{changes} at that threshold. A function with different
rules on different intervals is \ans{piecewise} --- exactly the last row of the
table in Part 2.
\end{notesbox}

\vspace{0.1in}

\boxguard
\begin{practicebox}
\begin{enumerate}[itemsep=8pt]
  \item Take differences until they are constant, then name the model type.

        \smallskip
        \renewcommand{\arraystretch}{1.4}
        \begin{center}
        \begin{tabular}{|c|c|c|c|c|c|}
        \hline
        $x$ & 0 & 1 & 2 & 3 & 4 \\
        \hline
        $y$ & 0 & 1 & 8 & 27 & 64 \\
        \hline
        1st & --- & \ans{1} & \ans{7} & \ans{19} & \ans{37} \\
        \hline
        2nd & --- & --- & \ans{6} & \ans{12} & \ans{18} \\
        \hline
        3rd & --- & --- & --- & \ans{6} & \ans{6} \\
        \hline
        \end{tabular}
        \end{center}

        Model type: \ans{cubic}

  \item A garden bed is built against a wall using 30 feet of edging on the
        other three sides. Build the area model and state its domain.

        \begin{work}
           2w + \ell &= 30 \\
                \ell &= 30 - 2w \\
                A(w) &= w(30 - 2w) \\
                     &= 30w - 2w^2
        \end{work}

        $A(w) = $ \ans{$30w - 2w^2$} \qquad Domain: \ans{0} $< w <$ \ans{15}
\end{enumerate}
\end{practicebox}

\end{document}
